% chapter01.tex — Basic Properties of Numbers (source: calcu5.pdf).
% Part divider and Disraeli epigraph live in main.tex, not here.

\spivakchapter{1}{Basic Properties of Numbers}

The title of this chapter expresses in a few words the mathematical knowledge
required to read this book. In fact, this short chapter is simply an
explanation of what is meant by the ``basic properties of numbers,'' all of
which---addition and multiplication, subtraction and division, solutions of
equations and inequalities, factoring and other algebraic manipulations---are
already familiar to us. Nevertheless, this chapter is not a review. Despite
the familiarity of the subject, the survey we are about to undertake will
probably seem quite novel; it does not aim to present an extended review of
old material, but to condense this knowledge into a few simple and obvious
properties of numbers. Some may even seem too obvious to mention, but a
surprising number of diverse and important facts turn out to be consequences
of the ones we shall emphasize.

Of the twelve properties which we shall study in this chapter, the first nine
are concerned with the fundamental operations of addition and multiplication.
For the moment we consider only addition: this operation is performed on a
pair of numbers---the sum $a+b$ exists for any two given numbers $a$ and $b$
(which may possibly be the same number, of course). It might seem reasonable
to regard addition as an operation which can be performed on several numbers
at once, and consider the sum $a_1 + \cdots + a_n$ of $n$ numbers
$a_1, \dots, a_n$ as a basic concept. It is more convenient, however, to
consider addition of pairs of numbers only, and to define other sums in terms
of sums of this type. For the sum of three numbers $a$, $b$, and $c$, this
may be done in two different ways. One can first add $b$ and $c$, obtaining
$b+c$, and then add $a$ to this number, obtaining $a+(b+c)$; or one can first
add $a$ and $b$, and then add the sum $a+b$ to $c$, obtaining $(a+b)+c$. Of
course, the two compound sums obtained are equal, and this fact is the very
first property we shall list:

\begin{quote}
\prop{1}\ \ If $a$, $b$, and $c$ are any numbers, then
\[ a + (b + c) = (a + b) + c. \]
\end{quote}

The statement of this property clearly renders a separate concept of the sum
of three numbers superfluous; we simply agree that $a+b+c$ denotes the number
$a+(b+c)=(a+b)+c$. Addition of four numbers requires similar, though slightly
more involved, considerations. The symbol $a+b+c+d$ is defined to mean
\begin{align*}
&(1) & ((a+b)+c)+d,\\
\text{or } &(2) & (a+(b+c))+d,\\
\text{or } &(3) & a+((b+c)+d),\\
\text{or } &(4) & a+(b+(c+d)),\\
\text{or } &(5) & (a+b)+(c+d).
\end{align*}
This definition is unambiguous since these numbers are all equal.
Fortunately, \emph{this} fact need not be listed separately, since it follows
from the property P1 already listed. For example, we know from P1 that
\[ (a+b)+c = a+(b+c), \]
and it follows immediately that (1) and (2) are equal. The equality of (2)
and (3) is a direct consequence of P1, although this may not be apparent at
first sight (one must let $b+c$ play the role of $b$ in P1, and $d$ the role
of $c$). The equalities $(3)=(4)=(5)$ are also simple to prove.

It is probably obvious that an appeal to P1 will also suffice to prove the
equality of the 14 possible ways of summing five numbers, but it may not be so
clear how we can reasonably arrange a proof that this is so without actually
listing these 14 sums. Such a procedure is feasible, but would soon cease to
be if we considered collections of six, seven, or more numbers; it would be
totally inadequate to prove the equality of all possible sums of an arbitrary
finite collection of numbers $a_1, \dots, a_n$. This fact may be taken for
granted, but for those who would like to worry about the proof (and it is
worth worrying about once) a reasonable approach is outlined in Problem~24.
Henceforth, we shall usually make a tacit appeal to the results of this
problem and write sums $a_1 + \cdots + a_n$ with a blithe disregard for the
arrangement of parentheses.

The number $0$ has one property so important that we list it next:

\begin{quote}
\prop{2}\ \ If $a$ is any number, then
\[ a + 0 = 0 + a = a. \]
\end{quote}

An important role is also played by $0$ in the third property of our list:

\begin{quote}
\prop{3}\ \ For every number $a$, there is a number $-a$ such that
\[ a + (-a) = (-a) + a = 0. \]
\end{quote}

Property P2 ought to represent a distinguishing characteristic of the number
$0$, and it is comforting to note that we are already in a position to prove
this. Indeed, if a number $x$ satisfies
\[ a + x = a \]
for any one number $a$, then $x=0$ (and consequently this equation also holds
for all numbers $a$). The proof of this assertion involves nothing more than
subtracting $a$ from both sides of the equation, in other words, adding $-a$
to both sides; as the following detailed proof shows, all three properties
P1--P3 must be used to justify this operation.
\begin{align*}
&\text{If} & a + x &= a,\\
&\text{then} & (-a) + (a+x) &= (-a) + a = 0;\\
&\text{hence} & ((-a)+a)+x &= 0;\\
&\text{hence} & 0 + x &= 0;\\
&\text{hence} & x &= 0.
\end{align*}

As we have just hinted, it is convenient to regard subtraction as an
operation derived from addition: we consider $a-b$ to be an abbreviation for
$a+(-b)$. It is then possible to find the solution of certain simple
equations by a series of steps (each justified by P1, P2, or P3) similar to
the ones just presented for the equation $a+x=a$. For example:
\begin{align*}
&\text{If} & x + 3 &= 5,\\
&\text{then} & (x+3)+(-3) &= 5 + (-3);\\
&\text{hence} & x + (3+(-3)) &= 5 - 3 = 2;\\
&\text{hence} & x + 0 &= 2;\\
&\text{hence} & x &= 2.
\end{align*}
Naturally, such elaborate solutions are of interest only until you become
convinced that they can always be supplied. In practice, it is usually just a
waste of time to solve an equation by indicating so explicitly the reliance
on properties P1, P2, and P3 (or any of the further properties we shall
list).

Only one other property of addition remains to be listed. When considering
the sums of three numbers $a$, $b$, and $c$, only two sums were mentioned:
$(a+b)+c$ and $a+(b+c)$. Actually, several other arrangements are obtained if
the order of $a$, $b$, and $c$ is changed. That these sums are all equal
depends on

\begin{quote}
\prop{4}\ \ If $a$ and $b$ are any numbers, then
\[ a + b = b + a. \]
\end{quote}

The statement of P4 is meant to emphasize that although the operation of
addition of pairs of numbers might conceivably depend on the order of the two
numbers, in fact it does not. It is helpful to remember that not all
operations are so well behaved. For example, subtraction does not have this
property: usually $a-b \neq b-a$. In passing we might ask just when $a-b$ does
equal $b-a$, and it is amusing to discover how powerless we are if we rely
only on properties P1--P4 to justify our manipulations. Algebra of the most
elementary variety shows that $a-b=b-a$ only when $a=b$. Nevertheless, it is
impossible to derive this fact from properties P1--P4; it is instructive to
examine the elementary algebra carefully and determine which step(s) cannot
be justified by P1--P4. We will indeed be able to justify all steps in detail
when a few more properties are listed. Oddly enough, however, the crucial
property involves multiplication.

The basic properties of multiplication are fortunately so similar to those
for addition that little comment will be needed; both the meaning and the
consequences should be clear. (As in elementary algebra, the product of $a$
and $b$ will be denoted by $a \cdot b$, or simply $ab$.)

\begin{quote}
\prop{5}\ \ If $a$, $b$, and $c$ are any numbers, then
\[ a \cdot (b \cdot c) = (a \cdot b) \cdot c. \]
\end{quote}

\begin{quote}
\prop{6}\ \ If $a$ is any number, then
\[ a \cdot 1 = 1 \cdot a = a. \]
\end{quote}

Moreover, $1 \neq 0$.

(The assertion that $1 \neq 0$ may seem a strange fact to list, but we have
to list it, because there is no way it could possibly be proved on the basis
of the other properties listed---these properties would all hold if there
were only one number, namely, $0$.)

\begin{quote}
\prop{7}\ \ For every number $a \neq 0$, there is a number $a^{-1}$ such that
\[ a \cdot a^{-1} = a^{-1} \cdot a = 1. \]
\end{quote}

\begin{quote}
\prop{8}\ \ If $a$ and $b$ are any numbers, then
\[ a \cdot b = b \cdot a. \]
\end{quote}

One detail which deserves emphasis is the appearance of the condition
$a \neq 0$ in P7. This condition is quite necessary; since $0 \cdot b = 0$
for all numbers $b$, there is no number $0^{-1}$ satisfying
$0 \cdot 0^{-1} = 1$. This restriction has an important consequence for
division. Just as subtraction was defined in terms of addition, so division
is defined in terms of multiplication: The symbol $a/b$ means $a \cdot b^{-1}$.
Since $0^{-1}$ is meaningless, $a/0$ is also meaningless---division by $0$ is
\emph{always} undefined.

Property P7 has two important consequences. If $a \cdot b = a \cdot c$, it
does not necessarily follow that $b = c$; for if $a = 0$, then both
$a \cdot b$ and $a \cdot c$ are $0$, no matter what $b$ and $c$ are. However,
if $a \neq 0$, then $b = c$; this can be deduced from P7 as follows:
\begin{align*}
&\text{If} & a \cdot b &= a \cdot c \text{ and } a \neq 0,\\
&\text{then} & a^{-1} \cdot (a \cdot b) &= a^{-1} \cdot (a \cdot c);\\
&\text{hence} & (a^{-1} \cdot a) \cdot b &= (a^{-1} \cdot a) \cdot c;\\
&\text{hence} & 1 \cdot b &= 1 \cdot c;\\
&\text{hence} & b &= c.
\end{align*}
It is also a consequence of P7 that if $a \cdot b = 0$, then either $a = 0$ or
$b = 0$. In fact,
\begin{align*}
&\text{if} & a \cdot b &= 0 \text{ and } a \neq 0,\\
&\text{then} & a^{-1} \cdot (a \cdot b) &= 0;\\
&\text{hence} & (a^{-1} \cdot a) \cdot b &= 0;\\
&\text{hence} & 1 \cdot b &= 0;\\
&\text{hence} & b &= 0.
\end{align*}
(It may happen that $a = 0$ and $b = 0$ are both true; this possibility is
not excluded when we say ``either $a = 0$ or $b = 0$''; in mathematics ``or''
is always used in the sense of ``one or the other, or both.'')

This latter consequence of P7 is constantly used in the solution of
equations. Suppose, for example, that a number $x$ is known to satisfy
\[ (x-1)(x-2) = 0. \]
Then it follows that either $x-1=0$ or $x-2=0$; hence $x=1$ or $x=2$.

On the basis of the eight properties listed so far it is still possible to
prove very little. Listing the next property, which combines the operations
of addition and multiplication, will alter this situation drastically.

\begin{quote}
\prop{9}\ \ If $a$, $b$, and $c$ are any numbers, then
\[ a \cdot (b + c) = a \cdot b + a \cdot c. \]
\end{quote}

(Notice that the equation $(b+c) \cdot a = b \cdot a + c \cdot a$ is also
true, by P8.)

As an example of the usefulness of P9 we will now determine just when
$a-b = b-a$:
\begin{align*}
&\text{If} & a - b &= b - a,\\
&\text{then} & (a-b)+b &= (b-a)+b = b+(b-a);\\
&\text{hence} & a &= b + b - a;\\
&\text{hence} & a + a &= (b+b-a)+a = b+b.\\
&\text{Consequently} & a \cdot (1+1) &= b \cdot (1+1),\\
&\text{and therefore} & a &= b.
\end{align*}

A second use of P9 is the justification of the assertion $a \cdot 0 = 0$
which we have already made, and even used in a proof on page~6 (can you find
where?). This fact was not listed as one of the basic properties, even though
no proof was offered when it was first mentioned. With P1--P8 alone a proof
was not possible, since the number $0$ appears only in P2 and P3, which
concern addition, while the assertion in question involves multiplication.
With P9 the proof is simple, though perhaps not obvious: We have
\begin{align*}
a \cdot 0 + a \cdot 0 &= a \cdot (0 + 0)\\
&= a \cdot 0;
\end{align*}
as we have already noted, this immediately implies (by adding $-(a \cdot 0)$
to both sides) that $a \cdot 0 = 0$.

A series of further consequences of P9 may help explain the somewhat
mysterious rule that the product of two negative numbers is positive. To
begin with, we will establish the more easily acceptable assertion that
$(-a) \cdot b = -(a \cdot b)$. To prove this, note that
\begin{align*}
(-a) \cdot b + a \cdot b &= [(-a)+a] \cdot b\\
&= 0 \cdot b\\
&= 0.
\end{align*}
It follows immediately (by adding $-(a \cdot b)$ to both sides) that
$(-a) \cdot b = -(a \cdot b)$. Now note that
\begin{align*}
(-a) \cdot (-b) + [-(a \cdot b)] &= (-a) \cdot (-b) + (-a) \cdot b\\
&= (-a) \cdot [(-b) + b]\\
&= (-a) \cdot 0\\
&= 0.
\end{align*}
Consequently, adding $(a \cdot b)$ to both sides, we obtain
\[ (-a) \cdot (-b) = a \cdot b. \]
The fact that the product of two negative numbers is positive is thus a
consequence of P1--P9. In other words, \emph{if we want P1 to P9 to be true,
the rule for the product of two negative numbers is forced upon us.}

The various consequences of P9 examined so far, although interesting and
important, do not really indicate the significance of P9; after all, we could
have listed each of these properties separately. Actually, P9 is the
justification for almost all algebraic manipulations. For example, although
we have shown how to solve the equation
\[ (x-1)(x-2) = 0, \]
we can hardly expect to be presented with an equation in this form. We are
more likely to be confronted with the equation
\[ x^2 - 3x + 2 = 0. \]
The ``factorization'' $x^2 - 3x + 2 = (x-1)(x-2)$ is really a triple use of P9:
\begin{align*}
(x-1) \cdot (x-2) &= x \cdot (x-2) + (-1) \cdot (x-2)\\
&= x \cdot x + x \cdot (-2) + (-1) \cdot x + (-1) \cdot (-2)\\
&= x^2 + x[(-2)+(-1)] + 2\\
&= x^2 - 3x + 2.
\end{align*}

A final illustration of the importance of P9 is the fact that this property
is actually used every time one multiplies arabic numerals. For example, the
calculation
\[
\begin{array}{r}
13 \\
\times 24 \\
\hline
52 \\
26 \\
\hline
312
\end{array}
\]
is a concise arrangement for the following equations:
\begin{align*}
13 \cdot 24 &= 13 \cdot (2 \cdot 10 + 4)\\
&= 13 \cdot 2 \cdot 10 + 13 \cdot 4\\
&= 26 \cdot 10 + 52.
\end{align*}
(Note that moving 26 to the left in the above calculation is the same as
writing $26 \cdot 10$.) The multiplication $13 \cdot 4 = 52$ uses P9 also:
\begin{align*}
13 \cdot 4 &= (1 \cdot 10 + 3) \cdot 4\\
&= 1 \cdot 10 \cdot 4 + 3 \cdot 4\\
&= 4 \cdot 10 + 12\\
&= 4 \cdot 10 + 1 \cdot 10 + 2\\
&= (4+1) \cdot 10 + 2\\
&= 5 \cdot 10 + 2\\
&= 52.
\end{align*}

The properties P1--P9 have descriptive names which are not essential to
remember, but which are often convenient for reference. We will take this
opportunity to list properties P1--P9 together and indicate the names by
which they are commonly designated.

\begin{center}
\renewcommand{\arraystretch}{1.35}
\begin{tabular}{@{}p{0.58\textwidth}p{0.38\textwidth}@{}}
\prop{1} \textit{(Associative law for addition)} &
  $a+(b+c)=(a+b)+c$.\\[0.35em]
\prop{2} \textit{(Existence of an additive identity)} &
  $a+0=0+a=a$.\\[0.35em]
\prop{3} \textit{(Existence of additive inverses)} &
  $a+(-a)=(-a)+a=0$.\\[0.35em]
\prop{4} \textit{(Commutative law for addition)} &
  $a+b=b+a$.\\[0.35em]
\prop{5} \textit{(Associative law for multiplication)} &
  $a\cdot(b\cdot c)=(a\cdot b)\cdot c$.\\[0.35em]
\prop{6} \textit{(Existence of a multiplicative identity)} &
  $a\cdot 1=1\cdot a=a$;\ $1\neq 0$.\\[0.35em]
\prop{7} \textit{(Existence of multiplicative inverses)} &
  $a\cdot a^{-1}=a^{-1}\cdot a=1$, for $a\neq 0$.\\[0.35em]
\prop{8} \textit{(Commutative law for multiplication)} &
  $a\cdot b=b\cdot a$.\\[0.35em]
\prop{9} \textit{(Distributive law)} &
  $a\cdot(b+c)=a\cdot b+a\cdot c$.
\end{tabular}
\end{center}

The three basic properties of numbers which remain to be listed are concerned
with inequalities. Although inequalities occur rarely in elementary
mathematics, they play a prominent role in calculus. The two notions of
inequality, $a<b$ ($a$ is less than $b$) and $a>b$ ($a$ is greater than $b$),
are intimately related: $a<b$ means the same as $b>a$ (thus $1<3$ and $3>1$
are merely two ways of writing the same assertion). The numbers $a$
satisfying $a>0$ are called \textbf{positive}, while those numbers $a$
satisfying $a<0$ are called \textbf{negative}. While positivity can thus be
defined in terms of $<$, it is possible to reverse the procedure: $a<b$ can
be defined to mean that $b-a$ is positive. In fact, it is convenient to
consider the collection of all positive numbers, denoted by $P$, as the basic
concept, and state all properties in terms of $P$:

\begin{quote}
\prop{10}\ \ \textit{(Trichotomy law)} For every number $a$, one and only
one of the following holds:
\begin{romans}
\item $a = 0$,
\item $a$ is in the collection $P$,
\item $-a$ is in the collection $P$.
\end{romans}
\end{quote}

\begin{quote}
\prop{11}\ \ \textit{(Closure under addition)} If $a$ and $b$ are in $P$,
then $a+b$ is in $P$.
\end{quote}

\begin{quote}
\prop{12}\ \ \textit{(Closure under multiplication)} If $a$ and $b$ are in
$P$, then $a \cdot b$ is in $P$.
\end{quote}

These three properties should be complemented with the following definitions:
\begin{align*}
a > b &\quad\text{if}\quad a - b \text{ is in } P;\\
a < b &\quad\text{if}\quad b > a;\\
a \geq b &\quad\text{if}\quad a > b \text{ or } a = b;\\
a \leq b &\quad\text{if}\quad a < b \text{ or } a = b.
\end{align*}
Note, in particular,\footnote{There is one slightly perplexing feature of
the symbols $\geq$ and $\leq$. The statements $1+1 \leq 3$ and $1+1 \leq 2$
are both true, even though we know that $\leq$ could be replaced by $<$ in
the first, and by $=$ in the second. This sort of thing is bound to occur
when $\leq$ is used with specific numbers; the usefulness of the symbol is
revealed by a statement like Theorem~1 where equality holds for some values
of $a$ and $b$, while inequality holds for other values.}
that $a>0$ if and only if $a$ is in $P$.

All the familiar facts about inequalities, however elementary they may seem,
are consequences of P10--P12. For example, if $a$ and $b$ are any two
numbers, then precisely one of the following holds:
\begin{romans}
\item $a - b = 0$,
\item $a - b$ is in the collection $P$,
\item $-(a-b) = b - a$ is in the collection $P$.
\end{romans}
Using the definitions just made, it follows that precisely one of the
following holds:
\begin{romans}
\item $a = b$,
\item $a > b$,
\item $b > a$.
\end{romans}

A slightly more interesting fact results from the following manipulations. If
$a<b$, so that $b-a$ is in $P$, then surely $(b+c)-(a+c)$ is in $P$; thus, if
$a<b$, then $a+c<b+c$. Similarly, suppose $a<b$ and $b<c$. Then
\begin{align*}
&b - a \text{ is in } P,\\
&\text{and}\quad c - b \text{ is in } P,\\
&\text{so}\quad c - a = (c-b)+(b-a) \text{ is in } P.
\end{align*}
This shows that if $a<b$ and $b<c$, then $a<c$. (The two inequalities $a<b$
and $b<c$ are usually written in the abbreviated form $a<b<c$, which has the
third inequality $a<c$ almost built in.)

The following assertion is somewhat less obvious: If $a<0$ and $b<0$, then
$ab>0$. The only difficulty presented by the proof is the unraveling of
definitions. The symbol $a<0$ means, by definition, $0>a$, which means
$0-a=-a$ is in $P$. Similarly $-b$ is in $P$, and consequently, by P12,
$(-a)(-b)=ab$ is in $P$. Thus $ab>0$.

The fact that $ab>0$ if $a>0$, $b>0$ and also if $a<0$, $b<0$ has one special
consequence: $a^2>0$ if $a\neq 0$. Thus squares of nonzero numbers are always
positive, and in particular we have proved a result which might have seemed
sufficiently elementary to be included in our list of properties: $1>0$ (since
$1=1^2$).

The fact that $-a>0$ if $a<0$ is the basis of a concept which will play an
extremely important role in this book. For any number $a$, we define the
\textbf{absolute value} $\abs{a}$ of $a$ as follows:
\[
\abs{a} =
\begin{cases}
a, & a \geq 0\\
-a, & a \leq 0.
\end{cases}
\]
Note that $\abs{a}$ is always positive, except when $a=0$. For example, we have
$\abs{-3}=3$, $\abs{7}=7$, $\abs{1+\sqrt{2}-\sqrt{3}} = 1+\sqrt{2}-\sqrt{3}$,
and $\abs{1+\sqrt{2}-\sqrt{10}} = \sqrt{10}-\sqrt{2}-1$. In general, the most
straightforward approach to any problem involving absolute values requires
treating several cases separately, since absolute values are defined by cases
to begin with. This approach may be used to prove the following very
important fact about absolute values.

\begin{theorem}
For all numbers $a$ and $b$, we have
\[ \abs{a+b} \leq \abs{a} + \abs{b}. \]
\end{theorem}

\begin{proof}
We will consider 4 cases:
\begin{steps}
\item $a \geq 0$,\ $b \geq 0$;
\item $a \geq 0$,\ $b \leq 0$;
\item $a \leq 0$,\ $b \geq 0$;
\item $a \leq 0$,\ $b \leq 0$.
\end{steps}
In case (1) we also have $a+b \geq 0$, and the theorem is obvious; in fact,
\[ \abs{a+b} = a+b = \abs{a}+\abs{b}, \]
so that in this case equality holds.

In case (4) we have $a+b \leq 0$, and again equality holds:
\[ \abs{a+b} = -(a+b) = -a+(-b) = \abs{a}+\abs{b}. \]

In case (2), when $a \geq 0$ and $b \leq 0$, we must prove that
\[ \abs{a+b} \leq a - b. \]
This case may therefore be divided into two subcases. If $a+b \geq 0$, then we
must prove that
\[ a+b \leq a-b, \]
i.e., $b \leq -b$, which is certainly true since $b \leq 0$ and hence
$-b \geq 0$. On the other hand, if $a+b \leq 0$, we must prove that
\[ -a-b \leq a-b, \]
i.e., $-a \leq a$, which is certainly true since $a \geq 0$ and hence
$-a \leq 0$.

Finally, note that case (3) may be disposed of with no additional work, by
applying case (2) with $a$ and $b$ interchanged.
\end{proof}

Although this method of treating absolute values (separate consideration of
various cases) is sometimes the only approach available, there are often
simpler methods which may be used. In fact, it is possible to give a much
shorter proof of Theorem~1; this proof is motivated by the observation that
\[ \abs{a} = \sqrt{a^2}. \]
(Here, and throughout the book, $\sqrt{x}$ denotes the \emph{positive} square
root of $x$; this symbol is defined only when $x \geq 0$.) We may now observe
that
\begin{align*}
(\abs{a+b})^2 &= (a+b)^2 = a^2 + 2ab + b^2\\
&\leq a^2 + 2\abs{a} \cdot \abs{b} + b^2\\
&= \abs{a}^2 + 2\abs{a} \cdot \abs{b} + \abs{b}^2\\
&= (\abs{a}+\abs{b})^2.
\end{align*}
From this we can conclude that $\abs{a+b} \leq \abs{a}+\abs{b}$ because
$x^2 < y^2$ implies $x < y$, provided that $x$ and $y$ are both nonnegative;
a proof of this fact is left to the reader (Problem~5).

One final observation may be made about the theorem we have just proved: a
close examination of either proof offered shows that
\[ \abs{a+b} = \abs{a}+\abs{b} \]
if $a$ and $b$ have the same sign (i.e., are both positive or both negative),
or if one of the two is $0$, while
\[ \abs{a+b} < \abs{a}+\abs{b} \]
if $a$ and $b$ are of opposite signs.

We will conclude this chapter with a subtle point, neglected until now, whose
inclusion is required in a conscientious survey of the properties of numbers.
After stating property P9, we proved that $a-b=b-a$ implies $a=b$. The proof
began by establishing that
\[ a \cdot (1+1) = b \cdot (1+1), \]
from which we concluded that $a=b$. This result is obtained from the equation
$a \cdot (1+1) = b \cdot (1+1)$ by dividing both sides by $1+1$. Division by
$0$ should be avoided scrupulously, and it must therefore be admitted that the
validity of the argument depends on knowing that $1+1 \neq 0$. Problem~25 is
designed to convince you that this fact cannot possibly be proved from
properties P1--P9 alone! Once P10, P11, and P12 are available, however, the
proof is very simple: We have already seen that $1>0$; it follows that
$1+1>0$, and in particular $1+1 \neq 0$.

This last demonstration has perhaps only strengthened your feeling that it is
absurd to bother proving such obvious facts, but an honest assessment of our
present situation will help justify serious consideration of such details. In
this chapter we have assumed that numbers are familiar objects, and that
P1--P12 are merely explicit statements of obvious, well-known properties of
numbers. It would be difficult, however, to justify this assumption. Although
one learns how to ``work with'' numbers in school, just what numbers are,
remains rather vague. A great deal of this book is devoted to elucidating the
concept of numbers, and by the end of the book we will have become quite well
acquainted with them. But it will be necessary to work with numbers
throughout the book. It is therefore reasonable to admit frankly that we do
not yet thoroughly understand numbers; we may still say that, in whatever way
numbers are finally defined, they should certainly have properties P1--P12.

Most of this chapter has been an attempt to present convincing evidence that
P1--P12 are indeed basic properties which we should assume in order to deduce
other familiar properties of numbers. Some of the problems (which indicate
the derivation of other facts about numbers from P1--P12) are offered as
further evidence. It is still a crucial question whether P1--P12 actually
account for \emph{all} properties of numbers. As a matter of fact, we shall
soon see that they do not. In the next chapter the deficiencies of properties
P1--P12 will become quite clear, but the proper means for correcting these
deficiencies is not so easily discovered. The crucial additional basic
property of numbers which we are seeking is profound and subtle, quite unlike
P1--P12. The discovery of this crucial property will require all the work of
Part~II of this book. In the remainder of Part~I we will begin to see why
some additional property is required; in order to investigate this we will
have to consider a little more carefully what we mean by ``numbers.''

\subsection*{PROBLEMS}

\pnum{1.} Prove the following:
\begin{romans}
\item If $ax = a$ for some number $a \neq 0$, then $x = 1$.
\item $x^2 - y^2 = (x-y)(x+y)$.
\item If $x^2 = y^2$, then $x = y$ or $x = -y$.
\item $x^3 - y^3 = (x-y)(x^2+xy+y^2)$.
\item $x^n - y^n = (x-y)(x^{n-1} + x^{n-2}y + \cdots + xy^{n-2} + y^{n-1})$.
\item $x^3 + y^3 = (x+y)(x^2 - xy + y^2)$. (There is a particularly easy way
  to do this, using (iv), and it will show you how to find a factorization
  for $x^n + y^n$ whenever $n$ is odd.)
\end{romans}

\pnum{2.} What is wrong with the following ``proof''? Let $x=y$. Then
\begin{align*}
x^2 &= xy,\\
x^2 - y^2 &= xy - y^2,\\
(x+y)(x-y) &= y(x-y),\\
x+y &= y,\\
2y &= y,\\
2 &= 1.
\end{align*}

\pnum{3.} Prove the following:
\begin{romans}
\item $\dfrac{a}{b} = \dfrac{ac}{bc}$, if $b, c \neq 0$.
\item $\dfrac{a}{b} + \dfrac{c}{d} = \dfrac{ad + bc}{bd}$, if $b, d \neq 0$.
\item $(ab)^{-1} = a^{-1}b^{-1}$, if $a, b \neq 0$. (To do this you must
  remember the defining property of $(ab)^{-1}$.)
\item $\dfrac{a}{b} \cdot \dfrac{c}{d} = \dfrac{ac}{db}$, if $b, d \neq 0$.
\item $\dfrac{a}{b} \bigg/ \dfrac{c}{d} = \dfrac{ad}{bc}$, if $b, c, d \neq 0$.
\item If $b, d \neq 0$, then $\dfrac{a}{b} = \dfrac{c}{d}$ if and only if
  $ad = bc$. Also determine when $\dfrac{a}{b} = \dfrac{b}{a}$.
\end{romans}

\pnum{4.} Find all numbers $x$ for which
\begin{romans}
\item $4 - x < 3 - 2x$.
\item $5 - x^2 < 8$.
\item $5 - x^2 < -2$.
\item $(x-1)(x-3) > 0$. (When is a product of two numbers positive?)
\item $x^2 - 2x + 2 > 0$.
\item $x^2 + x + 1 > 2$.
\item $x^2 - x + 10 > 16$.
\item $x^2 + x + 1 > 0$.
\item $(x-\pi)(x+5)(x-3) > 0$.
\item $(x-\sqrt[3]{2})(x-\sqrt{2}) > 0$.
\item $2^x < 8$.
\item $x + 3^x < 4$.
\item $\dfrac{1}{x} + \dfrac{1}{1-x} > 0$.
\item $\dfrac{x-1}{x+1} > 0$.
\end{romans}

\pnum{5.} Prove the following:
\begin{romans}
\item If $a < b$ and $c < d$, then $a + c < b + d$.
\item If $a < b$, then $-b < -a$.
\item If $a < b$ and $c > d$, then $a - c < b - d$.
\item If $a < b$ and $c > 0$, then $ac < bc$.
\item If $a < b$ and $c < 0$, then $ac > bc$.
\item If $a > 1$, then $a^2 > a$.
\item If $0 < a < 1$, then $a^2 < a$.
\item If $0 \leq a < b$ and $0 \leq c < d$, then $ac < bd$.
\item If $0 \leq a < b$, then $a^2 < b^2$. (Use (viii).)
\item If $a, b \geq 0$ and $a^2 < b^2$, then $a < b$. (Use (ix), backwards.)
\end{romans}

\pnum{6.}
\begin{letters}
\item Prove that if $0 \leq x < y$, then $x^n < y^n$, $n = 1, 2, 3, \dots$.
\item Prove that if $x < y$ and $n$ is odd, then $x^n < y^n$.
\item Prove that if $x^n = y^n$ and $n$ is odd, then $x = y$.
\item Prove that if $x^n = y^n$ and $n$ is even, then $x = y$ or $x = -y$.
\end{letters}

\pnum{7.} Prove that if $0 < a < b$, then
\[ a < \sqrt{ab} < \frac{a+b}{2} < b. \]
Notice that the inequality $\sqrt{ab} \leq (a+b)/2$ holds for all $a, b \geq 0$.
A generalization of this fact occurs in Problem~2-22.

\pnum{*8.} Although the basic properties of inequalities were stated in terms
of the collection $P$ of all positive numbers, and $<$ was defined in terms of
$P$, this procedure can be reversed. Suppose that P10--P12 are replaced by
\begin{quote}
\textbf{(P$'$10)}\ \ For any numbers $a$ and $b$ one, and only one, of the
following holds:
\begin{romans}
\item $a = b$,
\item $a < b$,
\item $b < a$.
\end{romans}

\textbf{(P$'$11)}\ \ For any numbers $a$, $b$, and $c$, if $a < b$ and
$b < c$, then $a < c$.

\textbf{(P$'$12)}\ \ For any numbers $a$, $b$, and $c$, if $a < b$, then
$a + c < b + c$.

\textbf{(P$'$13)}\ \ For any numbers $a$, $b$, and $c$, if $a < b$ and
$0 < c$, then $ac < bc$.
\end{quote}
Show that P10--P12 can then be deduced as theorems.

\pnum{9.} Express each of the following with at least one less pair of
absolute value signs.
\begin{romans}
\item $\abs{\sqrt{2} + \sqrt{3} - \sqrt{5} + \sqrt{7}}$.
\item $\abs{\abs{a+b} - \abs{a} - \abs{b}}$.
\item $\abs{\abs{a+b} + \abs{c} - \abs{a+b+c}}$.
\item $\abs{x^2 - 2xy + y^2}$.
\item $\abs{\abs{\sqrt{2} + \sqrt{3}} - \abs{\sqrt{5} - \sqrt{7}}}$.
\end{romans}

\pnum{10.} Express each of the following without absolute value signs,
treating various cases separately when necessary.
\begin{romans}
\item $\abs{a+b} - \abs{b}$.
\item $\abs{\abs{x}-1}$.
\item $\abs{x} - \abs{x^2}$.
\item $a - \abs{a-\abs{a}}$.
\end{romans}

\pnum{11.} Find all numbers $x$ for which
\begin{romans}
\item $\abs{x-3} = 8$.
\item $\abs{x-3} < 8$.
\item $\abs{x+4} < 2$.
\item $\abs{x-1} + \abs{x-2} > 1$.
\item $\abs{x-1} + \abs{x+1} < 2$.
\item $\abs{x-1} + \abs{x+1} < 1$.
\item $\abs{x-1} \cdot \abs{x+1} = 0$.
\item $\abs{x-1} \cdot \abs{x+2} = 3$.
\end{romans}

\pnum{12.} Prove the following:
\begin{romans}
\item $\abs{xy} = \abs{x} \cdot \abs{y}$.
\item $\abs*{\dfrac{1}{x}} = \dfrac{1}{\abs{x}}$, if $x \neq 0$. (The best
  way to do this is to remember what $\abs{x}^{-1}$ is.)
\item $\dfrac{\abs{x}}{\abs{y}} = \abs*{\dfrac{x}{y}}$, if $y \neq 0$.
\item $\abs{x-y} \leq \abs{x} + \abs{y}$. (Give a very short proof.)
\item $\abs{x} - \abs{y} \leq \abs{x-y}$. (A very short proof is possible, if
  you write things in the right way.)
\item $\abs{\abs{x}-\abs{y}} \leq \abs{x-y}$. (Why does this follow immediately
  from (v)?)
\item $\abs{x+y+z} \leq \abs{x}+\abs{y}+\abs{z}$. Indicate when equality
  holds, and prove your statement.
\end{romans}

\pnum{13.} The maximum of two numbers $x$ and $y$ is denoted by $\max(x,y)$.
Thus $\max(-1,3) = \max(3,3) = 3$ and $\max(-1,-4) = \max(-4,-1) = -1$. The
minimum of $x$ and $y$ is denoted by $\min(x,y)$. Prove that
\[
\max(x,y) = \frac{x + y + \abs{y-x}}{2},
\qquad
\min(x,y) = \frac{x + y - \abs{y-x}}{2}.
\]
Derive a formula for $\max(x,y,z)$ and $\min(x,y,z)$, using, for example
\[ \max(x,y,z) = \max(x,\max(y,z)). \]

\pnum{14.}
\begin{letters}
\item Prove that $\abs{a} = \abs{-a}$. (The trick is not to become confused
  by too many cases. First prove the statement for $a \geq 0$. Why is it then
  obvious for $a \leq 0$?)
\item Prove that $-b \leq a \leq b$ if and only if $\abs{a} \leq b$. In
  particular, it follows that $-\abs{a} \leq a \leq \abs{a}$.
\item Use this fact to give a new proof that $\abs{a+b} \leq \abs{a}+\abs{b}$.
\end{letters}

\pnum{*15.} Prove that if $x$ and $y$ are not both $0$, then
\[
x^2 + xy + y^2 > 0,
\qquad
x^4 + x^3y + x^2y^2 + xy^3 + y^4 > 0.
\]
\textit{Hint:} Use Problem~1.

\pnum{*16.}
\begin{letters}
\item Show that
\[
(x+y)^2 = x^2 + y^2 \text{ only when } x = 0 \text{ or } y = 0,
\]
\[
(x+y)^3 = x^3 + y^3 \text{ only when } x = 0 \text{ or } y = 0 \text{ or }
x = -y.
\]
\item Using the fact that
\[ x^2 + 2xy + y^2 = (x+y)^2 \geq 0, \]
show that $4x^2 + 6xy + 4y^2 > 0$ unless $x$ and $y$ are both $0$.
\item Use part (b) to find out when $(x+y)^4 = x^4 + y^4$.
\item Find out when $(x+y)^5 = x^5 + y^5$. \textit{Hint:} From the assumption
  $(x+y)^5 = x^5 + y^5$ you should be able to derive the equation
  $x^3 + 2x^2y + 2xy^2 + y^3 = 0$, if $xy \neq 0$. This implies that
  $(x+y)^3 = x^2y + xy^2 = xy(x+y)$.
\end{letters}
You should now be able to make a good guess as to when $(x+y)^n = x^n + y^n$;
the proof is contained in Problem~11-63.

\pnum{17.}
\begin{letters}
\item Find the smallest possible value of $2x^2 - 3x + 4$. \textit{Hint:}
  ``Complete the square,'' i.e., write $2x^2 - 3x + 4 = 2(x - 3/4)^2 + \text{?}$
\item Find the smallest possible value of $x^2 - 3x + 2y^2 + 4y + 2$.
\item Find the smallest possible value of $x^2 + 4xy + 5y^2 - 4x - 6y + 7$.
\end{letters}

\pnum{18.}
\begin{letters}
\item Suppose that $b^2 - 4c \geq 0$. Show that the numbers
\[ \frac{-b + \sqrt{b^2 - 4c}}{2}, \qquad \frac{-b - \sqrt{b^2 - 4c}}{2} \]
both satisfy the equation $x^2 + bx + c = 0$.
\item Suppose that $b^2 - 4c < 0$. Show that there are no numbers $x$
  satisfying $x^2 + bx + c = 0$; in fact, $x^2 + bx + c > 0$ for all $x$.
  \textit{Hint:} Complete the square.
\item Use this fact to give another proof that if $x$ and $y$ are not both
  $0$, then $x^2 + xy + y^2 > 0$.
\item For which numbers $a$ is it true that $x^2 + axy + y^2 > 0$ whenever
  $x$ and $y$ are not both $0$?
\item Find the smallest possible value of $x^2 + bx + c$ and of
  $ax^2 + bx + c$, for $a > 0$.
\end{letters}

\pnum{19.} The fact that $a^2 \geq 0$ for all numbers $a$, elementary as it
may seem, is nevertheless the fundamental idea upon which most important
inequalities are ultimately based. The great-granddaddy of all inequalities
is the \emph{Schwarz inequality}
\[
x_1 y_1 + x_2 y_2
\leq \sqrt{x_1^2 + x_2^2}\,\sqrt{y_1^2 + y_2^2}.
\]
(A more general form occurs in Problem~2-21.) The three proofs of the Schwarz
inequality outlined below have only one thing in common---their reliance on
the fact that $a^2 \geq 0$ for all $a$.

\begin{letters}
\item Prove that if $x_1 = \lambda y_1$ and $x_2 = \lambda y_2$ for some
  number $\lambda \geq 0$, then equality holds in the Schwarz inequality.
  Prove the same thing if $y_1 = y_2 = 0$. Now suppose that $y_1$ and $y_2$
  are not both $0$, and that there is no number $\lambda$ such that
  $x_1 = \lambda y_1$ and $x_2 = \lambda y_2$. Then
  \begin{align*}
  0 &< (\lambda y_1 - x_1)^2 + (\lambda y_2 - x_2)^2\\
    &= \lambda^2(y_1^2 + y_2^2) - 2\lambda(x_1 y_1 + x_2 y_2)
       + (x_1^2 + x_2^2).
  \end{align*}
  Using Problem~18, complete the proof of the Schwarz inequality.
\item Prove the Schwarz inequality by using $2xy \leq x^2 + y^2$ (how is this
  derived?) with
  \[
  x = \frac{x_i}{\sqrt{x_1^2 + x_2^2}},
  \qquad
  y = \frac{y_i}{\sqrt{y_1^2 + y_2^2}},
  \]
  first for $i = 1$ and then for $i = 2$.
\item Prove the Schwarz inequality by first proving that
  \[
  (x_1^2 + x_2^2)(y_1^2 + y_2^2)
  = (x_1 y_1 + x_2 y_2)^2 + (x_1 y_2 - x_2 y_1)^2.
  \]
\item Deduce, from each of these three proofs, that equality holds only when
  $y_1 = y_2 = 0$ or when there is a number $\lambda \geq 0$ such that
  $x_1 = \lambda y_1$ and $x_2 = \lambda y_2$.
\end{letters}

In our later work, three facts about inequalities will be crucial. Although
proofs will be supplied at the appropriate point in the text, a personal
assault on these problems is infinitely more enlightening than a perusal of a
completely worked-out proof. The statements of these propositions involve
some weird numbers, but their basic message is very simple: if $x$ is close
enough to $x_0$, and $y$ is close enough to $y_0$, then $x+y$ will be close to
$x_0 + y_0$, and $xy$ will be close to $x_0 y_0$, and $1/y$ will be close to
$1/y_0$. The symbol ``$\varepsilon$'' which appears in these propositions is
the fifth letter of the Greek alphabet (``epsilon''), and could just as well
be replaced by a less intimidating Roman letter; however, tradition has made
the use of $\varepsilon$ almost sacrosanct in the contexts to which these
theorems apply.

\pnum{20.} Prove that if
\[
\abs{x - x_0} < \frac{\varepsilon}{2}
\qquad\text{and}\qquad
\abs{y - y_0} < \frac{\varepsilon}{2},
\]
then
\[
\abs{(x+y) - (x_0+y_0)} < \varepsilon,
\qquad
\abs{(x-y) - (x_0-y_0)} < \varepsilon.
\]

\pnum{*21.} Prove that if
\[
\abs{x - x_0} < \min\left(\frac{\varepsilon}{2(\abs{y_0}+1)}, 1\right)
\qquad\text{and}\qquad
\abs{y - y_0} < \frac{\varepsilon}{2(\abs{x_0}+1)}
\]
then $\abs{xy - x_0 y_0} < \varepsilon$.
(The notation ``min'' was defined in Problem~13, but the formula provided by
that problem is irrelevant at the moment; the first inequality in the
hypothesis just means that
\[
\abs{x - x_0} < \frac{\varepsilon}{2(\abs{y_0}+1)}
\qquad\text{and}\qquad
\abs{x - x_0} < 1;
\]
at one point in the argument you will need the first inequality, and at
another point you will need the second. One more word of advice: since the
hypotheses only provide information about $x - x_0$ and $y - y_0$, it is
almost a foregone conclusion that the proof will depend upon writing
$xy - x_0 y_0$ in a way that involves $x - x_0$ and $y - y_0$.)

\pnum{*22.} Prove that if $y_0 \neq 0$ and
\[
\abs{y - y_0} < \min\left(\frac{\abs{y_0}}{2}, \frac{\varepsilon \abs{y_0}^2}{2}\right)
\]
then $y \neq 0$ and
\[
\abs*{\frac{1}{y} - \frac{1}{y_0}} < \varepsilon.
\]

\pnum{*23.} Replace the question marks in the following statement by
expressions involving $\varepsilon$, $x_0$, and $y_0$ so that the conclusion
will be true:
If $y_0 \neq 0$ and
\[
\abs{y - y_0} < \text{?} \qquad\text{and}\qquad \abs{x - x_0} < \text{?}
\]
then $y \neq 0$ and
\[
\abs*{\frac{x}{y} - \frac{x_0}{y_0}} < \varepsilon.
\]
This problem is trivial in the sense that its solution follows from Problems
21 and 22 with almost no work at all (notice that $x/y = x \cdot 1/y$). The
crucial point is not to become confused; decide which of the two problems
should be used first, and don't panic if your answer looks unlikely.

\pnum{*24.} This problem shows that the actual placement of parentheses in a
sum is irrelevant. The proofs involve ``mathematical induction''; if you are
not familiar with such proofs, but still want to tackle this problem, it can
be saved until after Chapter~2, where proofs by induction are explained.

Let us agree, for definiteness, that $a_1 + \cdots + a_n$ will denote
\[
a_1 + (a_2 + (a_3 + \cdots + (a_{n-2} + (a_{n-1} + a_n))) \cdots).
\]
Thus $a_1 + a_2 + a_3$ denotes $a_1 + (a_2 + a_3)$, and $a_1 + a_2 + a_3 + a_4$
denotes $a_1 + (a_2 + (a_3 + a_4))$, etc.

\begin{letters}
\item Prove that
\[
(a_1 + \cdots + a_k) + a_{k+1} = a_1 + \cdots + a_{k+1}.
\]
\textit{Hint:} Use induction on $k$.
\item Prove that if $n \geq k$, then
\[
(a_1 + \cdots + a_k) + (a_{k+1} + \cdots + a_n) = a_1 + \cdots + a_n.
\]
\textit{Hint:} Use part (a) to give a proof by induction on $k$.
\item Let $s(a_1, \dots, a_k)$ be some sum formed from $a_1, \dots, a_k$.
  Show that
\[ s(a_1, \dots, a_k) = a_1 + \cdots + a_k. \]
\textit{Hint:} There must be two sums $s'(a_1, \dots, a_l)$ and
$s''(a_{l+1}, \dots, a_k)$ such that
\[ s(a_1, \dots, a_k) = s'(a_1, \dots, a_l) + s''(a_{l+1}, \dots, a_k). \]
\end{letters}

\pnum{25.} Suppose that we interpret ``number'' to mean either $0$ or $1$, and
$+$ and $\cdot$ to be the operations defined by the following two tables.
\begin{center}
\begin{tabular}{c|cc}
$+$ & $0$ & $1$ \\
\hline
$0$ & $0$ & $1$ \\
$1$ & $1$ & $0$
\end{tabular}
\qquad\qquad
\begin{tabular}{c|cc}
$\cdot$ & $0$ & $1$ \\
\hline
$0$ & $0$ & $0$ \\
$1$ & $0$ & $1$
\end{tabular}
\end{center}
Check that properties P1--P9 all hold, even though $1 + 1 = 0$.
